\section{Further research questions}
\label{sec:research}

The completed results suggest specific next steps. We distinguish a
request for a sharper proof from a conjecture about the mathematics;
the absence of a reduction in a finite search is never used as an
independence conjecture by itself.

\subsection{Gaussian identities and proof compression}

\paragraph{1. Compress the $S_4$ certificate.}
Find a short derivation of Theorem~\ref{thm:s4} from a small collection
of named functional identities. The present certificate is a complete
proof, but its $911$ standard rows obscure which intermediate values
are essential. A useful computational formulation is sparse exact
elimination with penalties on depth, color changes, and coefficient
height. Every proposed compression should be replayed by the independent
verifier. This is a concrete optimization problem with an already
certified target.

\paragraph{2. Determine the higher even-index mixed reductions.}
Apply the convergent involution to
$\Li_{1,2m}(i,1)$ and to other low-depth seeds at weights $2m+1$.
For $m=3,4$, ask whether $S_{2m}$ can be expressed using the same-point
doubles $\operatorname{Im}\Li_{a,2m+1-a}(i,1)$ and an explicitly
specified algebra of lower-weight products. The phrase ``explicitly
specified'' is essential: changing the allowed basis changes the
problem. A positive answer requires a rational certificate; a negative
answer in a chosen formal relation system only describes that system.

\paragraph{3. Organize the octahedral relations by representations.}
The puncture set $\{0,\infty,\pm1,\pm i\}$ has a rich symmetry group.
Its action on iterated-integral words suggests decomposing the formal
weight space before elimination. Determine which symmetry component
contains the $S_4$ target and why one additional convergent seed suffices
for this particular calculation. The aim is a theorem about the relation
module, with motivic and numerical period claims stated separately.

\paragraph{4. Formalize the finite certificate.}
The standard-library verifier has a deliberately small interface:
word shuffles, stuffles, sign distribution, conjugation, and one
M\"obius substitution. These are suitable for translation into a proof
assistant. The mathematical task is to connect the formal relation
schemas to the analytic definitions, including the single-divergence
regularization. Checking a JSON residual alone would formalize only
the linear algebra, not that analytic connection.

\subsection{Algebraic bases and higher ladders}

\paragraph{5. Prove the plastic tetralogarithm.}
The adjacent manuscript formula, in reduced coefficients, is
\begin{align}
0={}&-71\Li_4(q)-15\Li_4(q^2)+39\Li_4(q^3)-17\Li_4(q^5)
       +16\Li_4(q^7)-\Li_4(q^{14})\notag\\
 &\quad+51\zeta(4)+\frac{335}{4}\lambda^4
                 -42\zeta(2)\lambda^2,
 \qquad q^3+q^2=1,\quad\lambda=\log q.
 \label{eq:open-plastic-four}
\end{align}
This is an inherited, numerically supported target, not a new
conjecture introduced by this article. The low-weight proofs suggest
searching for a specialization of a tetralogarithm functional equation
whose arguments close after introducing $q^7$ and $q^{14}$. A symbol or
cobracket calculation can identify candidates, but the remaining
constant and branch terms still need proof.

\paragraph{6. Move from individual complements to product relations.}
Classify, or effectively enumerate at bounded support and exponent
height, relations
$\prod_j(1-r^j)^{n_j}=r^N$ in a specified number field.
The theorem in this article shows that individual pure complements
have very limited redundancy. Higher ladders therefore invite a
systematic study of product relations and of the additional units
they introduce. Exact ideal valuations, units, and cyclotomic
factorizations offer a natural algebraic starting point.

\paragraph{7. Refine the degree count.}
The exact formula for $N(D)$ reduces its error term beyond $D^2/2$ to
an explicit sum involving $e(\lfloor D/d\rfloor+1)$ and
$e(\lfloor D/d\rfloor+2)$. Determine its secondary asymptotic behavior,
including any $D\log D$ contribution and arithmetic fluctuations.
The formula permits large exact experiments without polynomial
factorization. Such experiments should precede any proposed secondary
constant; this article does not conjecture an unevidenced one.

\paragraph{8. Allow two independent generators.}
For multiplicatively independent positive algebraic numbers $r,s$,
study solutions to $x+y=1$ in
$\{\pm r^as^b:a,b\in\mathbb Z\}$. General unit-equation theory gives
the surrounding framework, but a sharp, explicit analogue of the
cyclic classification would require additional arithmetic input.
Begin with fields and generators already appearing in the manuscript,
and state the equivalence under a change of generators carefully.

\subsection{Uniform asymptotics and gamma geometry}

\paragraph{9. Connect the proportional and critical harmonic regimes.}
Construct an error estimate uniform as $\alpha=p/h$ grows from a
positive constant to $\log h+O(1)$. The present rate satisfies
$I(\alpha)\sim e^{-\alpha}/2$, while the preceding report has a
critical limit $\exp(-e^{-c}/2)$ for $\alpha=\log h+c$.
These formulas match formally, but a connecting theorem needs moving
saddles, parameter-dependent curvature, and quantitative tail bounds.
It should recover both the present $D(\alpha,a)$ correction and the
earlier critical correction in their overlapping range.

\paragraph{10. Let both reflected gamma indices grow.}
For $m/n\to\theta>0$, the integrand can be written
\[
 \exp\{n[\log(\log\Gamma(x))
              +\theta\log(\log\Gamma(1-x))]\}.
\]
The endpoint mechanisms of the fixed-$m$ theorem then change.
Locate the dominant point or points, establish uniqueness or a precise
bifurcation statement, and derive a uniform Laplace expansion. At
$\theta=1$, symmetry forces the phase to be symmetric about $1/2$;
it does not by itself prove that $1/2$ is the global maximum.

\paragraph{11. Determine the optimally truncated signed remainder.}
The unsigned bound is sufficient for certified approximation, but it
does not identify the leading exponentially small remainder when
$K$ lies near the least term. Develop a contour or inverse-function
description that determines its sign, scale, and dependence on the
fractional part of the optimal truncation index. The already proved
late-coefficient law supplies the candidate least-term location;
turning it into a remainder theorem is a separate analytic step.

\paragraph{12. Classify early residue zeros and sign changes.}
The first coefficient is positive, and the late sign is
$(-1)^{k+m-1}$. Early behavior depends on $m$: even the second
nonzero coefficient changes sign between $m=1$ and $m=2$.
Determine whether any early coefficient vanishes and obtain explicit
bounds for the onset of the eventual sign law as a function of $m$.
An exact identity in the formal zeta--gamma polynomial algebra should
be distinguished from a numerical zero after substituting Euler's
constant and zeta values.

\paragraph{13. Develop weighted inverse-gamma moments.}
Replace $B(x)^m$ by an analytic weight $W(x)$ with a specified order
of vanishing at zero. The residue calculation suggests
\[
 A_k(W)=\frac1k[x^{k-1}]\Gamma(1+x)^kW(x),
\]
and the same Cauchy-circle method should yield a late-coefficient
expansion when the transformed amplitude does not vanish at the
critical point. Precisely identify the hypotheses under which this
holds, including endpoint integrability, zeros at the saddle, and
weights with their own singularities. Reflected integer powers are
the fully proved case in this article.

These questions have different proof requirements. Exact word
identities and algebraic searches are natural targets for small
certificates. Uniform asymptotics require analytic remainder estimates.
Statements about numerical period independence require arithmetic
ideas beyond the reduction and numerical methods used here.
